% ============================================================
% Model: CuO$_2$ 平面的对称性与布拉维格子
% ============================================================

\section{CuO$_2$ 平面的对称性与布拉维格子}
\label{sec:cuo2-symmetry}

\begin{tipbox}[关键词标签]
布拉维格子 · 点群 · 空间群 · CuO$_2$ 平面 · 高温超导 · 4mm 点群 · 结构相变
\end{tipbox}

% --------------------------------------------------------
% 1. 模型定义框
% --------------------------------------------------------
\begin{modelbox}[模型：CuO$_2$ 平面的晶体学]
\textbf{物理情境}：CuO$_2$ 平面是铜基高温超导体（如 YBCO、BSCCO）的核心结构单元。Cu 原子位于正方格子的格点上，O 原子位于 Cu--Cu 连线的中点（或略有偏移）。理解其对称性是分析电子结构和超导机理的基础。

\textbf{CuO$_2$ 平面结构}：
\begin{itemize}[nosep]
    \item Cu 原子：位于 $(0,0)$、$(a,0)$、$(0,a)$、$(a,a)$ 等（正方布拉维格子）
    \item O 原子：位于 $(a/2,0)$、$(0,a/2)$ 及其平移等价位置
    \item 每个 Cu 与四个 O 配位，形成 CuO$_4$ 平面正方配位
\end{itemize}

\textbf{关键对称操作}：
\begin{itemize}[nosep]
    \item 四重旋转轴 $C_4$（垂直于平面）
    \item 镜面 $\sigma_v$（沿对角线和对边中点连线）
    \item 镜面 $\sigma_d$（沿 Cu--O 键方向和对角线）
\end{itemize}
\end{modelbox}

% --------------------------------------------------------
% 2. 关键公式框
% --------------------------------------------------------
\begin{formulabox}[核心公式]
\begin{enumerate}[label=\textbf{(\arabic*)}]
    \item \textbf{CuO$_2$ 的基元}：
    \[
    \text{Cu 在 }(0,0),\quad\text{O 在 }\Bigl(\frac{a}{2},0\Bigr),\Bigl(0,\frac{a}{2}\Bigr)
    \]
    \texteq{每个原胞含 1 个 Cu 和 2 个 O，共 3 个原子}

    \item \textbf{正方格子的倒格矢}：
    \[
    \mathbf{b}_1=\frac{2\pi}{a}(1,0),\quad\mathbf{b}_2=\frac{2\pi}{a}(0,1)
    \]
    \texteq{倒格子也是正方格子，第一布里渊区为正方形}

    \item \textbf{布里渊区高对称点}：
    \[
    \Gamma=(0,0),\quad X=\Bigl(\frac{\pi}{a},0\Bigr),\quad M=\Bigl(\frac{\pi}{a},\frac{\pi}{a}\Bigr)
    \]

    \item \textbf{点群 $4mm$ 的生成元}：
    \[
    C_4\;(z\text{ 轴}),\quad\sigma_v\;(xz),\quad\sigma_v'\;(yz)
    \]
    \texteq{$4mm$ 是正方晶系的点群，含 8 个对称操作}

    \item \textbf{正交畸变后的点群}：
    \[
    4mm\xrightarrow{a\neq b}mm2\;(C_{2v})
    \]
    \texteq{$C_4$ 轴被破坏，仅剩 $C_2$ 和两个垂直镜面}
\end{enumerate}
\end{formulabox}

% --------------------------------------------------------
% 3. 训练点框
% --------------------------------------------------------
\begin{drillbox}[训练点与考法变体]
\trainingpoint{1} \textbf{布拉维格子判定}：给定 CuO$_2$ 平面原子排布，画出原胞，确定布拉维格子类型（正方格子，非简单正方）。

\trainingpoint{2} \textbf{点群分析}：列出 CuO$_2$ 平面的所有对称操作，证明其点群为 $4mm$（国际符号）或 $C_{4v}$（Schoenflies 符号）。

\trainingpoint{3} \textbf{畸变后的对称性降低}：当正方畸变为正交（$a\neq b$），点群降为 $mm2$（$C_{2v}$），说明哪些对称操作丢失。

\trainingpoint{4} \textbf{空间群}：CuO$_2$ 平面是二维结构，空间群为 $p4mm$（No.~11）。写出其 Seitz 符号或对称操作集合。

\trainingpoint{5} \textbf{与铜基超导关联}：CuO$_2$ 平面中的 $d_{x^2-y^2}$ 轨道与 O 的 $p$ 轨道杂化形成反键带，费米面附近的电子态决定超导性质。结构畸变影响电子结构（如 Jahn-Teller 效应）。
\end{drillbox}

% --------------------------------------------------------
% 4. 解题技巧框
% --------------------------------------------------------
\begin{tipbox}[快速识别与解题技巧]
\begin{itemize}
    \item \examnote{识别标志}：题干出现``CuO$_2$''''高温超导''''正方格子''''点群 4mm''，立即定位本模型。
    \item \examnote{正方 vs 正交判定}：若 $a=b$ 且 $\gamma=90°$，为四方/正方；若 $a\neq b$ 且 $\gamma=90°$，为正交。注意题目给出的晶格常数是否相等。
    \item \examnote{点群元素计数}：$4mm$ 含 8 个元素：$E, C_4, C_4^3, C_2, \sigma_v, \sigma_v', \sigma_d, \sigma_d'$。注意 $C_4^2=C_2$，且镜面有两套（沿轴和对角线）。
    \item \examnote{2D 布拉维格子速查}：二维只有 5 种布拉维格子：斜方、长方（含中心）、正方、六角。
    \item \examnote{畸变导致能隙}：正方到正交的畸变会破坏 $C_4$ 对称性，使原本简并的能带（如 $d_{xz}$ 和 $d_{yz}$）分裂，可能在费米面打开赝能隙。
\end{itemize}
\end{tipbox}

% --------------------------------------------------------
% 5. 易错点框
% --------------------------------------------------------
\begin{errorbox}{常见失误与陷阱}
\begin{itemize}
    \item \textbf{CuO$_2$ 不是简单正方}：虽然 Cu 原子形成简单正方格子，但 O 原子位于边心位置，整个结构是``正方格子 + 基元''，而非简单正方布拉维格子（布拉维格子只看等价点）。
    \item \textbf{O 原子位置}：O 原子只在 $x$ 轴和 $y$ 轴方向的 Cu--Cu 中点，不在对角线中点。如果误认为 O 在 $(a/2,a/2)$，则对称性分析完全错误。
    \item \textbf{2D 与 3D 点群}：2D 点群只有 10 种（$1,2,3,4,6,m,2mm,3m,4mm,6mm$），3D 点群有 32 种。CuO$_2$ 平面分析用 2D 点群。
    \item \textbf{4mm vs 4/mmm}：$4mm$ 是 2D 点群，无反演中心；$4/mmm$ 是 3D 点群（四方晶系全对称点群），含反演。CuO$_2$ 平面只有 $4mm$。
    \item \textbf{畸变方向}：正交畸变通常是沿 $a$ 轴伸长、沿 $b$ 轴压缩（或反之），不是沿对角线。这导致 $C_4$ 被破坏，但 $C_2$（沿 $z$ 轴）保留。
\end{itemize}
\end{errorbox}

% --------------------------------------------------------
% 6. 物理图像与关联模型
% --------------------------------------------------------
\begin{insightbox}[物理图像与模型关联]
\textbf{物理图像}：

CuO$_2$ 平面就像一个棋盘：
\begin{itemize}[nosep]
    \item Cu 在棋盘格的交点上，O 在格子的边中点。
    \item 如果你把棋盘旋转 90°，看起来完全一样——这就是 $C_4$ 对称性。
    \item 如果你沿水平或垂直中线对折，也完全重合——这就是 $\sigma_v$ 镜面。
    \item 如果你沿对角线对折，仍然重合——这就是 $\sigma_d$ 镜面。
    \item 但如果你把棋盘拉长（变成矩形），90° 旋转不再对称——$C_4$ 丢失，只剩 $C_2$（180° 旋转）。
\end{itemize}

\textbf{与相似模型的对比}：
\begin{center}
\begin{tabular}{llll}
\toprule
\textbf{对比维度} & \textbf{CuO$_2$ 平面} & \textbf{石墨烯} & \textbf{La$_2$CuO$_4$} \\
\midrule
布拉维格子 & 正方 & 六角 & 四方 \\
点群 & $4mm$ & $6mm$ & $4/mmm$（3D） \\
配位 & 平面四配位 & 平面三配位 & 八面体六配位 \\
电子性质 & 金属/关联绝缘体 & 半金属 & 反铁磁绝缘体 \\
\bottomrule
\end{tabular}
\end{center}

\textbf{推广方向}：
\begin{itemize}[nosep]
    \item Jahn-Teller 效应：Cu$^{2+}$（$d^9$）在平面四配位中发生 $e_g$ 轨道分裂，导致晶格畸变
    \item 条纹相（stripe phase）：CuO$_2$ 平面中自旋和电荷的周期性调制，破坏平移对称性
    \item 双层 CuO$_2$：Bi-2212 等含双层 CuO$_2$，层间耦合导致能带劈裂（bonding/antibonding）
\end{itemize}
\end{insightbox}

% --------------------------------------------------------
% 7. 推导附录
% --------------------------------------------------------
\begin{derivationbox}[推导细节（自我检验用）]
\textbf{Step 1：CuO$_2$ 平面的对称操作}

以 Cu 在 $(0,0)$，O 在 $(a/2,0)$ 和 $(0,a/2)$ 为例：

\begin{itemize}
    \item $E$：恒等操作，显然保持结构不变。
    \item $C_4$（$z$ 轴旋转 90°）：$(x,y)\to(-y,x)$。
    $(a/2,0)\to(0,a/2)$，O 位置互换，结构不变。
    \item $C_4^3$（270° 旋转）：类似 $C_4$。
    \item $C_2$（180° 旋转）：$(x,y)\to(-x,-y)$，O 位置映射到等价位置。
    \item $\sigma_v$（$xz$ 平面反射）：$(x,y)\to(x,-y)$。
    $(a/2,0)\to(a/2,0)$（O 不动），$(0,a/2)\to(0,-a/2)$（等价 O 位置）。
    \item $\sigma_v'$（$yz$ 平面反射）：类似。
    \item $\sigma_d$（对角线 $y=x$ 反射）：$(x,y)\to(y,x)$。
    $(a/2,0)\to(0,a/2)$，O 位置互换。
    \item $\sigma_d'$（对角线 $y=-x$ 反射）：类似。
\end{itemize}

共 8 个对称操作，构成点群 $4mm$（$C_{4v}$）。

\textbf{Step 2：正交畸变后的对称性}

畸变后 $a\neq b$（设 $a>b$）：
\begin{itemize}
    \item $C_4$：$(a/2,0)\to(0,a/2)$，但 $a\neq b$ 使 $(0,a/2)$ 与 $(b/2,0)$ 不等价，$C_4$ 不保持结构。
    \item $C_2$：$(x,y)\to(-x,-y)$，仍保持（因为 $a$ 和 $b$ 的符号变化不影响距离）。
    \item $\sigma_v$（$xz$）：$(x,y)\to(x,-y)$，仍保持（$a$ 不变，$b$ 变号不影响）。
    \item $\sigma_v'$（$yz$）：类似，保持。
    \item $\sigma_d$：$(x,y)\to(y,x)$，$(a/2,0)\to(0,a/2)$，但 $a\neq b$ 时不等价，不保持。
\end{itemize}

剩余操作：$E, C_2, \sigma_v, \sigma_v'$，构成点群 $mm2$（$C_{2v}$）。
\end{derivationbox}

\vspace{1em}
